\documentclass[a4paper,10pt]{article}

\usepackage{amscd,amssymb,amsmath,graphicx,verbatim,hyperref}
\usepackage{textcomp}
\usepackage[OT1,T1]{fontenc}

\usepackage[all,graph,curve]{xypic}
\usepackage{pdfpages}
\makeatletter
 \def\setlabelmargin#1{\labelmargin@=#1\relax }
\makeatother

%\newcommand{\qed}{\nobreak \ifvmode \relax \else
%      \ifdim\lastskip<1.5em \hskip-\lastskip
%      \hskip1.5em plus0em minus0.5em \fi \nobreak
%      \vrule height0.75em width0.5em depth0.25em\fi}

\newcounter{Definitioncount}
\setcounter{Definitioncount}{0}

\newtheorem{theorem}{Theorem}
\newtheorem{lemma}[theorem]{Lemma}
\newtheorem{proposition}[theorem]{Proposition}
\newtheorem{corollary}[theorem]{Corollary}
\newenvironment{example}[1][Example]{\begin{trivlist}
\item[\hskip \labelsep {\bfseries #1}] }{\end{trivlist}}
\newenvironment{proof}[1][Proof]{\begin{trivlist}
\item[\hskip \labelsep {\bfseries #1}]}{\end{trivlist}}
\newenvironment{definition}[1][Definition]{\begin{trivlist}
\item[\hskip \labelsep {\bfseries #1}] \refstepcounter{Definitioncount} \textbf{\arabic{Definitioncount}} }{\end{trivlist}}
\newenvironment{remark}[1][Remark]{\begin{trivlist}
\item[\hskip \labelsep {\bfseries #1}]}{\end{trivlist}}

%Extra stuff
\newcommand{\fc}{\mathscr{C}}
\newcommand{\bm}{{\cal{M}}}
\newcommand{\bid}{{\cal{I}}}
\newcommand{\bc}{{\cal{C}}}
\newcommand{\ba}{{\cal{A}}}
\newcommand{\bk}{{\cal{K}}}
\newcommand{\by}{{\cal{Y}}}
\newcommand{\bx}{{\cal{X}}}
\newcommand{\bB}{{\cal{B}}}
\newcommand{\duof}{{\cal{F}}}
\newcommand{\oop}{\operatorname{op}}
\newcommand{\op}{^{\oop}}
\newcommand{\pps}{\operatorname{ps}}
\newcommand{\ps}{^{\pps}}
\newcommand{\lowps}{_{\pps}}
\newcommand{\ttract}{\operatorname{tract}}
\newcommand{\tract}{_{\ttract}}
\newcommand{\rrep}{\operatorname{rep}}
\newcommand{\rep}{_{\rrep}}
\newcommand{\bicat}{\mathrm{Bicat}(\bc\op,\bm\op)\op}
\newcommand{\ev}{\mathcal{V}}
\newcommand{\lom}{\ell{om}}
\newcommand{\rom}{r{om}}
\newcommand{\warmon}{\,\Box\,}
\numberwithin{equation}{section}


\begin{document}

\thispagestyle{empty}

\textbf{{\Large Some diagrammatic monoidalities}}
\newline
\newline

\noindent Only nodes:

$$\xymatrix{
& A \otimes (B \otimes C) & (B\otimes C) \otimes A & \\
(A\otimes B)\otimes C & & & B \otimes (C \otimes A) \\
& (B \otimes A) \otimes C & B \otimes (A \otimes C) &
}$$
\newline
\newline

\noindent With arrows and all frivolities:

$$\xymatrix{
& A \otimes (B \otimes C) \ar[r]^-{\gamma} & (B\otimes C) \otimes A \ar[rd]^-{\alpha} & \\
(A\otimes B)\otimes C \ar[ru]^-{\alpha} \ar[rd]_-{\gamma \otimes 1} & & & B \otimes (C \otimes A) \\
& (B \otimes A) \otimes C \ar[r]_-{\alpha} & B \otimes (A \otimes C) \ar[ru]_-{1 \otimes \gamma} &
}$$
\newline
\newline

\noindent A pentagon:

$$\xymatrix{
& (A \otimes (B \otimes C)) \otimes D \ar[rd]^-{\phantom{A}\alpha_{A,B\otimes C,D}}  & \\
((A \otimes B )\otimes C) \otimes D \ar[ru]^-{\alpha_{A,B,C}\otimes 1_{D}\phantom{A}} \ar[d]_-{\alpha_{A\otimes B,C,D}} & & A \otimes ((B \otimes C) \otimes D ) \ar[d]^-{1_{A}\otimes\alpha_{B,C,D}} \\
(A \otimes B) \otimes (C \otimes D) \ar[rr]_-{\alpha_{A,B,C\otimes D }} & & A \otimes (B \otimes (C \otimes D)) 
}$$
\newline

\noindent You'll notice that I've used some phantom characters in the above diagram's label code.
Try removing them to see what happens.
Things like this are great little quick and dirty tricks to fix some common positional and layout problems.

\newpage

\thispagestyle{empty}


\noindent A triangle:

$$\xymatrix{
(A\otimes I)\otimes B \ar[rr]^-{\alpha_{A,I,B}} \ar[rd]_-{\rho_{A}\otimes 1_{B}} && A\otimes (I\otimes B) \ar[ld]^-{1_{A}\otimes \lambda_{B}} \\
& A\otimes B &
}$$
\newline
\newline

\noindent Here is a version of the pentagon diagram that's been made too big to fit on the page:

$$\xymatrix{
& & (A \otimes (B \otimes C)) \otimes D \ar[rrdd]^-{\phantom{A}\alpha_{A,B\otimes C,D}}  & & \\
& & & & \\
((A \otimes B )\otimes C) \otimes D \ar[rruu]^-{\alpha_{A,B,C}\otimes 1_{D}\phantom{A}} \ar[rdd]_-{\alpha_{A\otimes B,C,D}} & & & & A \otimes ((B \otimes C) \otimes D ) \ar[ldd]^-{1_{A}\otimes\alpha_{B,C,D}} \\
& & & & \\
& (A \otimes B) \otimes (C \otimes D) \ar[rr]_-{\alpha_{A,B,C\otimes D }} && A \otimes (B \otimes (C \otimes D)) &
}$$
\newline
\newline


\noindent Short of redesigning it so that it does fit, one may simply scale the image by encapsulating the entire diagram's code in a scale-box with a scale factor between $0$ and $1$.
Here I've used $0.75$:

$$\scalebox{0.75}{\xymatrix{
& & (A \otimes (B \otimes C)) \otimes D \ar[rrdd]^-{\phantom{A}\alpha_{A,B\otimes C,D}}  & & \\
& & & & \\
((A \otimes B )\otimes C) \otimes D \ar[rruu]^-{\alpha_{A,B,C}\otimes 1_{D}\phantom{A}} \ar[rdd]_-{\alpha_{A\otimes B,C,D}} & & & & A \otimes ((B \otimes C) \otimes D ) \ar[ldd]^-{1_{A}\otimes\alpha_{B,C,D}} \\
& & & & \\
& (A \otimes B) \otimes (C \otimes D) \ar[rr]_-{\alpha_{A,B,C\otimes D }} && A \otimes (B \otimes (C \otimes D)) &
}}$$
\newline

\noindent In most situations this is not ideal (including this one in my opinion) and it is better to choose a different style of diagram.
In fact it would be better to use xygraph or xy instead of xymatrix due to not being bound to the rigid matrix-entry-like positioning (i.e. in xy you can have nodes overlapping each other and, while you can force this in xymatrix, it is pretty hopeless).


\end{document}